\section{Ausgangslage: in welchem Kanal lebt $\theta=2/9$?}
Der Koide-Winkel $\theta=2/9$ ist empirisch auf sieben Stellen bestätigt, aber nicht aus der
Geometrie hergeleitet. Er ist selbst kein Ansatz, sondern eine \emph{gefundene} Größe: erst die
Hilbertraum-Übersetzung von FFGFT (Dok.~230/231/232) und ihre Fortführung bis zu den Massen
(Dok.~282) machte sichtbar, dass die drei Lepton-Massen das Spektrum eines hermiteschen
Z$_3$-Zirkulanten auf der $\mathbb{C}^3$-Faser sind, in dessen Eigenwerten $\theta=2/9$ als
Ergebnis der Diagonalisierung erscheint. Genau dieser Zirkulant liefert die Amplituden $a_k$
unten. Zur empirischen Einordnung des Winkels -- $2/9$ als rationale Adresse eines über einen
Referenzpunkt fixierten Wertes, nicht als Polmassen-Exaktheitsanspruch (P42) -- siehe Dok.~292;
für den hier gesuchten dynamischen Ort ist diese Feinunterscheidung ohne Belang, da die
Orbit-Auflösung weit über dem $10^{-7}$-Niveau liegt. Die vorangehende Analyse (Dok.~290) trennt
am Massenoperator einen
\emph{Betragskanal} (die Eigenwerte, die Massen) von einem \emph{Phasenkanal} (Mischung,
Allpass, Holonomie) und zeigt: ein Betragsspektrum legt die Phase nur bis auf einen
Allpass-Faktor fest. Die Frage dieses Dokuments ist entsprechend: In welchem Kanal sitzt
$\theta=2/9$ -- und welcher Art Größe ist er, verglichen mit einem erzwungenen Anker wie der
Lichtgeschwindigkeit $c$? Die Rechnungen dazu liegen in \texttt{allpass\_theta.py},
\texttt{allpass\_carrier\_selection.py} und \texttt{forced\_vs\_contingent.py} (numpy-only,
Seed 20780458).

\section{Die Betrags-Seite ist counter-blind}
Mit den vorzeichenrichtigen Amplituden $a_k(\theta)=1+\sqrt2\cos(\theta+2\pi k/3)$ ($k=0,1,2$)
ist der symmetrische Koide-Invariant $Q=(\sum_k a_k^2)/(\sum_k a_k)^2$ exakt und
\emph{$\theta$-unabhängig}. Denn die Summe $S=\sum_k a_k=3+\sqrt2\sum_k\cos(\theta+2\pi k/3)=3$
(drei um $120^\circ$ versetzte Kosinus summieren zu null), und die Quadratsumme
$T=\sum_k a_k^2=3+0+2\cdot\tfrac32=6$ (mit $\sum_k\cos^2=\tfrac32$). Damit
\[
Q(\theta)=\frac{T}{S^2}=\frac{6}{9}=\frac{2}{3}\qquad\text{für jedes }\theta.
\]
$Q$ läuft für positive Massen über $[1/3,1]$; der Wert $2/9=0{,}2222$ liegt \emph{unter} diesem
Bereich. $2/9$ ist somit kein Grenzwert, kein Extremum und kein Fixpunkt des symmetrischen
Funktionals -- im Betragskanal ist $2/9$ vollständig unsichtbar. Das ist die algebraische Seite
der \enquote{Counter-Blindheit} des $Z_3$-Kerns: die Orbit-Glieder $\{1/9,2/9,4/9\}$ sind für
jede symmetrische, zählende Größe ununterscheidbar.

\section{Der Allpass: drei Kanäle}
Der Ort von $\theta$ lässt sich am Allpass sichtbar machen. Ein $Z_3$-symmetrisches
Blaschke-Produkt $B_\theta(z)=\prod_{k}(z^{-1}-\bar a_k)/(1-a_k z^{-1})$ mit Polen
$a_k=r\,e^{i(\theta+2\pi k/3)}$ (hier $r=0{,}5$) zerlegt die Phaseninformation in drei Größen,
von denen zwei gegen $\theta$ blind sind und nur eine ihn trägt. Numerisch (für die
Orbit-Winkel $\theta\in\{1/9,2/9,4/9\}$):
\begin{itemize}
\item \textbf{Betrag} $|B_\theta|=1$ (bis $\sim10^{-15}$) für alle $\theta$ -- der Betragskanal
  ist blind.
\item \textbf{Windungszahl} $=3$ exakt für alle $\theta$ -- die \emph{flach-topologische}
  ganze Zahl ist blind. Das ist die numerische Seite der Aussage, dass $2/(9\pi)$ irrational
  und $\theta=2/9$ daher kein flacher topologischer Invariant ist: die ganzzahlige Windung
  trennt $2/9$ nicht von $1/9,4/9$.
\item \textbf{Stetiges Phasenprofil} $\arg B_\theta(\omega)$ -- verschieden für
  $1/9,2/9,4/9$ (am Referenzpunkt $\omega=0$: $0{,}093$, $0{,}171$, $0{,}249$; auch die
  Gruppenlaufzeit unterscheidet sie). Nur das \emph{stetige, dynamische} Phasenprofil trägt
  $\theta$.
\end{itemize}
Auch ein Energie-/Betrags-Funktional wie $\int|B|^2\,d\omega/2\pi$ ist für alle drei identisch.
Die Distinktion $2/9$ gegen $1/9,4/9$ lebt somit weder im Betrag, noch in der ganzzahligen
Windung, noch in einem symmetrischen Funktional -- allein im dynamischen Phasenprofil.

\section{Ein illustrativer Auswahlmechanismus}
Was ein Selektor leisten müsste, zeigt eine dynamische Phasen-Referenz. Die drei Orbit-Glieder
haben Allpass-Phasen $\psi=(0{,}093,\,0{,}171,\,0{,}249)$. Ein Träger (\emph{Carrier}), der eine
Holonomie $\Phi$ akkumuliert und auf die nächste Orbit-Phase einrastet, wählt genau einen
Repräsentanten: $2/9$ wird ausgewählt genau dann, wenn $\Phi$ im $2/9$-Basin
$[0{,}132,\,0{,}210]$ liegt (Breite $0{,}078$). Eine \emph{zufällige} Holonomie über den
Orbit-Bereich trifft $2/9$ mit Wahrscheinlichkeit $\approx1/3$ -- ohne Grund bevorzugt nichts
den Wert $2/9$.

Die Form dieses Mechanismus ist die des offenen Problems: die Auswahl ist \emph{kontingent} auf
eine externe dynamische Phase (eine Holonomie); $2/9$ hat ein eigenes Basin wie jedes Glied,
und Betrag wie Windung geben ihm keinen Vorzug. Ein dynamischer Selektor müsste diese
Holonomie \emph{erzwungen} liefern, aus einem inneren Grund der liefernden Dynamik, nicht durch
eine von Hand gewählte Umbenennung. Dies ist derselbe Inhalt wie die Brechungsbedingung $(\delta,\chi)\approx(0{,}24,
\pi/2)$: $\chi$ \emph{ist} die Holonomie-Phase, $\delta$ ihre freie Tiefe. (Illustration der
Mechanismus-Form; keine Herleitung des Wertes $2/9$, keine Modellierung eines konkreten
Trägers.)

\section{Erzwungen oder kontingent: $\theta^*(\delta)$ und der Kontrast zum Photon}
Ob $2/9$ \emph{erzwungen} ist wie $c$, entscheidet das Variationsverhalten. Der $\delta=0$-Anteil
des Brechungsfunktionals hat die geschlossene Form (bis $10^{-15}$ bestätigt)
\[
E_3(\theta,0)=-\tfrac12+\tfrac{1}{\sqrt2}\cos 3\theta,\qquad
\frac{dE_3}{d\theta}\Big|_{0}=-\frac{3}{\sqrt2}\sin 3\theta,
\]
mit Extrema an den $Z_3$-Punkten ($\sin3\theta=0$). Koppelt man das zweite Triplett $3'$ mit
Stärke $\delta$ und Phase $\chi$ an, wandern die inneren Extrema vom $Z_3$-Punkt weg (es sind
zwei; verfolgt wird der untere, orbitnahe Zweig $\theta^*$). Auf dem Zweig $\chi=\pi/2$ ist die
Verschiebung glatt und monoton (numerisch,
\texttt{forced\_vs\_contingent.py}): $\theta^*$ steigt von $0{,}099$ ($\delta=0{,}10$) über
$0{,}223$ ($\delta=0{,}24$) bis $0{,}501$ ($\delta=1{,}0$); $d\theta^*/d\delta$ ist überall
positiv (Median über das $\delta$-Gitter $0{,}31$, lokal bei $\delta^*\approx0{,}24$ etwa
$0{,}80$). Der erreichte Bereich $[0{,}05,\,0{,}52]$ enthält $2/9$ im Inneren, unausgezeichnet.
Das $2/9$-Basin im $\theta$-Raum -- nächstes Orbit-Glied, $[1/6,\,1/3]$, Breite $1/6$ --
entspricht über die lokale Steigung einem $\delta$-Fenster von
$\Delta\delta\approx0{,}167/0{,}80\approx0{,}21$. $\delta$ ist weich, nicht feinabgestimmt;
$2/9$ wird herbeigezogen, es sitzt an keiner Struktur.

Die Lichtgeschwindigkeit ist von anderem Typ. In $ds^2=-c^2dt^2+dx^2+dy^2+dz^2$ ist $c$ ein
metrischer Koeffizient, kein Parameter eines Funktionals; $ds^2=0$ ist der Lichtkegel, und
masselos bedeutet $v=c$ exakt -- die Sättigung der Ungleichung $v\le c$. Es gibt kein
$c^*(\lambda)$; $dc/d\lambda\equiv0$, weil $c$ eine Konstante der Geometrie ist, kein Extremum.
Dies spiegelt die physikalische Herkunft: die Masselosigkeit des Photons ist
\emph{eichsymmetrie-geschützt} ($m=0$ exakt, daraus $v=c$), ein \emph{Rand}-/Fixpunktwert.
$2/9$ dagegen ist der irreduzibel $Z_3$-\emph{brechende} Winkel im \emph{Inneren} des Orbits --
symmetrie-brechend, nicht -geschützt, und daher kontingent.

\section{Selbst- und schief-adjungierte Holonomien}
Ist der Selektor eine Holonomie-Phase $\chi$, so ist zwischen zwei Sektoren zu unterscheiden. Im
Brechungsterm $\delta\cos(2\varphi_k+\chi)$ ($\varphi_k=\theta+2\pi k/3$) trennt $\chi$ eine
\emph{gerade} von einer \emph{ungeraden} Kopplung: $\chi=0$ (und $\chi=\pi$) gibt
$\cos(2\varphi_k)$, eine gerade, \emph{selbst­adjungierte}, reelle Kopplung; $\chi=\pi/2$ gibt
$\cos(2\varphi_k+\pi/2)=-\sin(2\varphi_k)$, eine ungerade, \emph{schief-adjungierte},
imaginäre Kopplung.

Die beiden Sektoren haben verschiedene natürliche Werte. Die \emph{selbst­adjungierten}
Holonomien -- flache $Z_3$-Wilson-Linien ($\chi=2\pi n/3$) mal die zwei Spinstrukturen pro
Zyklus (Versatz $0$ oder $\pi$) -- sind die Vielfachen von $\pi/3$,
$\chi\in\{0,\pi/3,2\pi/3,\pi,4\pi/3,5\pi/3\}$; $\pi/2$ ist keiner davon. Der
\emph{schief-adjungierte} Sektor dagegen hat seinen natürlichen Wert bei $\chi=\pi/2$: das ist
die $i$-Phase, $e^{i\pi/2}=i=\sqrt{-1}$, gegenüber der reellen $e^{i\pi}=-1$. Der Übergang von
reeller zu schief-adjungierter Kopplung \emph{halbiert} die Phase $\pi\to\pi/2$ -- er zieht die
Wurzel aus $-1$.

Entscheidend ist, welcher Sektor den Selektor stellt. Die Eliminationskette (Dok.~282) hat alle
statischen, selbst­adjungierten Kandidaten ausgeschieden; der einzige Überlebende ist die
dynamische \emph{schief-adjungierte} Phase. Damit ist der typ-konsistente Wert $\chi=\pi/2$ --
die $i$-Phase des schief-adjungierten Sektors --, nicht ein Vielfaches von $\pi/3$. Und
$\chi=\pi/2$ erreicht $\theta^*=2/9$ bei $\delta\approx0{,}24$ (numerisch,
\texttt{holonomie\_landschaft.py}). Der schief-adjungierte Typ lässt dabei zunächst beide
$i$-Werte zu, $\chi=\pi/2$ und $\chi=3\pi/2$ ($e^{\pm i\pi/2}=\pm i$); die Wahl des Zweigs ist
eine Orientierungswahl (Umlaufsinn des Zyklus). Die beiden Zweige sind jedoch nicht äquivalent:
der Spiegel-Zweig $\chi=3\pi/2$ besitzt bei $\delta\approx0{,}24$ kein inneres Extremum nahe
$2/9$ (Extrema bei $0{,}83$ und $1{,}87$) -- nur der $+i$-Zweig erreicht den Orbit-Bereich.

$\chi=\pi/2$ ist damit kein willkürlicher oder topologisch heimatloser Wert, sondern der von der
Eliminationskette geforderte schief-adjungierte $i$-Wert -- ohne angenommene $Z_4$ und ohne
Zahlenspiel. Präzise gehört \enquote{schief-adjungiert} dabei zum \emph{Generator}, nicht zur
endlichen Holonomie: diese ist ein unitäres Gruppenelement $\exp(K)$ mit schief-adjungiertem
$K$ ($K^\dagger=-K$); der schief-adjungierte Charakter sitzt in der erzeugenden Kopplung, nicht
in der endlichen Holonomie-Matrix selbst. Die Tiefe $\delta$ bleibt jedoch \emph{frei}: der
Sektor legt die Phase $\chi$ fest, nicht die Tiefe. Der offene Punkt ist damit geschärft --
nicht ein beliebiger Winkel, sondern ob eine schief-adjungierte Dynamik gerade $\chi=\pi/2$
trägt und darüber die Tiefe $\delta$ aus einem inneren Grund \emph{dieser Dynamik} festlegt
(rekursionsintern ist nach Dok.~275 keine Auszeichnung möglich).

\section{Die Tiefe $\delta$ als zweite Harmonische}
Die Massen-Amplitude ist eine Fourier-Reihe im Faserwinkel,
\[
a(\varphi)=\underbrace{1}_{\text{Mittelwert}}+\underbrace{\sqrt2\cos\varphi}_{\text{1.\ Harmonische}}
+\underbrace{\delta\cos(2\varphi+\chi)}_{\text{2.\ Harmonische}},\qquad \varphi=\theta+2\pi k/3.
\]
Die erste Harmonische mit reeller Amplitude $c_1=\sqrt2$ legt den Koide-Invarianten $Q=2/3$ fest
(sie ist der Betrag $r=\sqrt2$). Die zweite Harmonische hat die komplexe Amplitude
$c_2=\delta\,e^{i\chi}$; ihr \emph{Betrag} ist $\delta=|c_2|$ und ihre \emph{Phase} ist
$\chi=\arg c_2$. Für $\chi=\pi/2$ wird $c_2=i\delta$ rein imaginär -- die zweite Harmonische ist
ein reiner Sinus, die schief-adjungierte Kopplung des vorigen Abschnitts.

Damit sind $\delta$ und $\chi$ Betrag und Phase \emph{derselben} Größe $c_2$, gehören aber zu
verschiedenen Sektoren: $\chi=\arg c_2$ liegt im Phasen-Sektor (von der Eliminationskette auf
$\pi/2$ fixiert), $\delta=|c_2|$ im \emph{Betrags}-Sektor. Die Tiefe $\delta$ ist also kein
Phasen-, sondern ein Amplituden-Freiheitsgrad -- die Amplitude des ersten Obertons, physikalisch
die Stärke, mit der das zweite Triplett $3'$ (die Verdopplung $6=3\oplus3'$) in das primäre $3$
einmischt.

Das ist derselbe Sektor, der bereits $|c_1|=\sqrt2$ festlegt. Wie die erste harmonische Amplitude
aus $Q=2/3$ folgt, müsste die zweite, $|c_2|=\delta$, aus dem Oberton-Gehalt der
massenerzeugenden fraktalen Wellenform folgen (der Struktur $D_f=3-\xi$, der
$3\oplus3'$-Verdopplung). Das Verhältnis $\delta/\sqrt2\approx0{,}17$ -- die zweite Harmonische ist
rund $17\,\%$ der ersten -- ist eine reine \emph{Form}-Größe, die an der Wellenform hängt, nicht
an der Phase. $\delta$ ist somit kein freier Abstimmparameter, sondern ein noch nicht
ausgerechneter Betrags-Freiheitsgrad.

Damit trennt sich die Frage nach $2/9$ endgültig in zwei Sektoren: die \emph{Phase} $\chi=\pi/2$
ist schief-adjungiert und fixiert; der \emph{Betrag} $\delta$ ist die zweite harmonische Amplitude
und gehört -- offen -- in denselben Betrags-Sektor wie $\sqrt2$.

\section{$\delta$ als golden gedämpfte Amplitude}
Der Oberton-Gehalt lässt sich schärfer fassen: die zweite harmonische Amplitude ist eine
\emph{Golden-Dämpfung}. Der Wert, der $\theta^*=2/9$ bei $\chi=\pi/2$ exakt liefert, ist
$\delta^*=0{,}2389$; der Golden-Attraktor $1/\varphi^3=0{,}2361$ liegt knapp darunter. Weil
$1/\varphi^3=\varphi^{-3}$ eine \emph{negative} Golden-Potenz (ein Zerfall) ist, ist $\delta$
eine Dämpfung, keine Verstärkung -- $\delta^*<1$ liegt eindeutig im Dämpfungs-Regime.

Der Exponent $3$ ist nicht frei, sondern $=D$. Eine einzelne, dekompaktifizierte
$3'$-Rekursion gäbe $1/\varphi^1=0{,}618$ (um $159\,\%$ daneben); erst die Dämpfung über
$D=3$ Dimensionen (die fraktale/räumliche Dimension $D_f=3-\xi$, das \emph{volle} kompakte
$T^4$) gibt $1/\varphi^3$. Der passende Exponent $3$ ist damit selbst ein Argument für das
volle kompakte $T^4$ und gegen die Dekompaktifizierung auf eine einzelne Rekursion.

Die $\xi$-Dämpfungs-Rekursion $r(k{+}1)=r(k)(1-\xi)$, $r(0)=D_f/3$ (Dok.~275) passiert
$1/\varphi$ bei der Tiefe $k^*=\ln\varphi/\xi\approx3609$ und $1/\varphi^3$ bei $3\cdot3609
\approx10827$ (auf $0{,}011\,\%$ genau). $\delta^*$ wird rund $89$ \emph{Einheiten} davor
erreicht -- gemeint ist eine Differenz zweier reeller Log-Tiefen im dekompaktifizierten,
kontinuierlichen Skalenfluss ($1/\varphi^3$ bei $10826{,}2$, $\delta^*$ bei $10737{,}7$;
Differenz $88{,}5$), keine ganzzahlige Schrittzählung. So ist $\delta^*$ ein um rund $89$
Log-Einheiten untergedämpftes $1/\varphi^3$, ein Überschuss von $+1{,}187\,\%$ über dem
Attraktor -- beide Angaben sind über $\ln(\delta^*\varphi^{3})/\xi$ ein und dieselbe Zahl,
keine zwei unabhängigen Befunde --, konsistent mit einer Annäherung von oben.

Ob dieser Überschuss die fraktale Massenkorrektur $K_{\text{frak}}=1-100\xi$ (Dok.~133) ist,
lässt sich beim gegenwärtigen Präzisionsstand \emph{nicht entscheiden}. Deren präzise Form
$(D_f^{\text{eff}}/3)^{D_f^{\text{eff}}/2}$ mit $D_f^{\text{eff}}\approx2{,}973$ ergibt einen Faktor
$\approx100$ ($1{,}333\,\%$), der gemessene Überschuss beträgt $\approx89\,\xi$ ($1{,}187\,\%$) --
ein Abstand von nur $0{,}147\,\%$. Weil $\delta^*$ über $\theta^*=2/9$ an die $\tau$-limitierte
Winkelbestimmung gebunden ist ($\mathrm{d}\theta^*/\mathrm{d}\delta\approx0{,}31$), trägt es selbst
eine Toleranz von $\approx0{,}135\,\%$; die beiden Kandidaten $89\xi$ und $100\xi$ liegen also nur
rund $1\sigma$ auseinander und sind mit heutiger Präzision \emph{nicht unterscheidbar}. Eine
Ablehnung von $K_{\text{frak}}$ auf $0{,}15\,\%$ wäre schärfer als die eigene Bestimmungstoleranz von
$\delta$ und damit unzulässig -- ganz analog zur $N_0$-Bewertung (Dok.~292: $100/\varphi^2$ liegt bei
$1{,}6\sigma$ und ist \emph{zulässig}, nicht widerlegt). $K_{\text{frak}}$ ist überdies strukturell
motiviert: die zweite Harmonische und $K_{\text{frak}}$ entspringen \emph{derselben} fraktalen
Korrektur $g=\eta(1+\xi f)$ (Dok.~286), von der $K_{\text{frak}}$ die RG-Ebene ist (Dok.~133). Damit
ist $K_{\text{frak}}$ ein \emph{zulässiger, strukturell motivierter Kandidat} für den Überschuss,
nicht ausgeschlossen.

Zwei Grenzen bleiben. Erstens: zulässig heißt nicht hergeleitet -- die Toleranz entfernt die falsche
Ablehnung, singelt $K_{\text{frak}}$ im Band aber nicht aus (dort liegen $89$, $100$ und alles
dazwischen); die Auszeichnung vor einem beliebigen Bandwert trägt allein das Struktur-Argument, nicht
die Zahl. Zweitens: der gemeinsame $g=\eta(1+\xi f)$-Ursprung ist in Dok.~286 für den \emph{Phasen}-Kanal
(Holonomie) gezeigt; der Transfer auf den \emph{Betrag} der zweiten Harmonischen ($\delta$) ist
plausibel, aber noch nicht gerechnet. Unabhängig davon fixiert Dok.~275 die Rekursionstiefe ohnehin
nicht: $1/\varphi^3$ ist dort ein \emph{Durchgangspunkt}, kein Fixpunkt (der einzige Fixpunkt der
Dämpfung ist $0$), und zu \emph{jedem} Ziel $c\in(0,1)$ existiert eine Tiefe $k^*(c)$ mit identischem
Parameterstatus. Dass $\delta^*$ nahe $1/\varphi^3$ liegt, ist daher \emph{kein} Schließen der
Rekursion; der Feinwert bleibt ein zulässiger, aber nicht hergeleiteter Kandidat.

Damit ist der Betrags-Sektor so weit geklärt, wie die Struktur trägt: $\delta$ ist die golden
gedämpfte zweite Harmonische, $1/\varphi^3$ über $D=3$ (volles $T^4$), mit einem kleinen
Überschuss von oben. Form, Ordnung, Exponent und Richtung stehen fest; der exakte $\delta$-Feinwert
-- der Überschuss über $1/\varphi^3$ -- ist mit $K_{\text{frak}}=1-100\xi$ innerhalb der
$\theta$-Toleranz \emph{verträglich} und strukturell motiviert (Dok.~286, gemeinsamer
$g=\eta(1+\xi f)$-Ursprung), aber nicht als Betragswert hergeleitet und bleibt insoweit offen
(\texttt{delta\_golden\_daempfung.py}).

\section{Was die Nicht-Schließung praktisch bedeutet: ein offener Kopplungskanal}
Der bisherige Befund trennt zwei Sektoren mit gegensätzlichem Verhalten. Der Massen-Sektor ist
\emph{geschlossen}: die Rekursion $R(\Phi^*)=\Phi^*$ (Dok.~203) besitzt einen nicht-trivialen
Fixpunkt, und weil die Lepton-Massen an ihm sitzen, sind sie und der Koide-Wert $Q=2/3$ exakt und
starr.\footnote{\enquote{Exakt} bezieht sich hier auf die geometrische Fixpunkt-Ebene: der
dimensionslose Invariant $Q=2/3$ ist starr und wird von den Daten auf $-0{,}9\,\sigma$
bestätigt (Dok.~292). Die \emph{Absolut}massen bzw.\ das Zahlenpaar $(\sqrt2,\ 2/9)$ sind davon
zu unterscheiden: auf der gemessenen Polmassen-Ebene ist das Paar nicht gemeinsam exakt (Dok.~292,
$\sim450\,\sigma$ im $m_\mu/m_e$-Kanal), weshalb $2/9$ per P42 als rationale Adresse eines über
einen Referenzpunkt fixierten Winkels gilt, nicht als Polmassen-Exaktheitsanspruch. Für den hier
diskutierten dynamischen Ort ist das ohne Belang -- die Orbit-Auflösung ($\sim0{,}1$) liegt weit
über dem $10^{-7}$-Unterschied.}
Dort gibt es nichts zu verstellen -- kein freier Parameter, kein Angriffspunkt für eine
äußere Beeinflussung.

Der Spielraum sitzt allein im offenen Sektor, der Phase. Die $\xi$-Dämpfungs-Rekursion schließt
nie (Durchgangspunkt, einziger Fixpunkt $0$, Dok.~275); in der Sprache der Rückkopplung ist sie
eine offene, gedämpfte Schleife mit Verstärkung kleiner eins, kein starr eingerasteter Fixpunkt.
Eine solche offene, gedämpfte Schleife kann -- anders als ein geschlossener Fixpunkt -- von außen
beeinflusst werden: durch Kopplung an einen zweiten Oszillator, durch Injektions- oder
Mode-Locking (Dok.~286). Praktisch heißt die Nicht-Schließung daher: es existiert ein Kanal, über
den ein zweites, dynamisches System die offene Phase einrasten kann. Der geschlossene
Massen-Sektor bietet diesen Kanal nicht.

Der Kanal ist jedoch eng und strukturiert, kein freier Parameter. Der Fangbereich einer solchen
Kopplung ist von der Ordnung $\xi$ (Dok.~286) -- winzig; die $2/9$-Arnold-Zunge ist rund
zwölfmal schmaler als $1/3$. Die Kopplung benachbarter Skalenebenen im kontinuierlichen Fluss ist
real, aber betragsmäßig von der Ordnung $10^{-13}$ und destabilisiert die geschlossenen
Vorhersagen nicht. Und die Wirkung ist gerichtet: sie läuft von der Dynamik in die Phase, nicht in
den radialen Massen-Sektor, der starr bleibt (der radial/angulare Schnitt aus Dok.~282).

Der eine Ort, an dem dieser Kanal genutzt würde, ist ein dynamischer schief-adjungierter Generator
außerhalb der statischen FFGFT-Geometrie -- die Fortsetzung der Eliminationskette (Dok.~283/286).
Rastet ein solcher Generator die offene Phase gerade auf $\chi=\pi/2$ und damit $\theta=2/9$ ein,
erhielte $2/9$ den dynamischen, geschützten Status, den die folgende Bilanz beschreibt. Das ist
ausdrücklich ein Kandidat, kein erreichtes Resultat: FFGFTs eigene Rekursion $R$ ist rein reell und
trägt keinen solchen Generator; ob eine äußere Dynamik ihn liefert, bleibt offen.

Entscheidend ist die Grenze dieses Spielraums. \enquote{Raum für gegenseitige Beeinflussung}
bedeutet nicht, dass sich die Theorie beliebig biegen ließe. Der Raum ist beschränkt: der radiale
Sektor bleibt starr, die Phase ist eng eingeschnürt (schief-adjungiert, $\chi=\pi/2$, transzendent
in $2\pi$, $\xi$-schmaler Fangbereich), und die Rekursionstiefe fixiert keinen Wert (Dok.~275). Es
ist strukturierter Spielraum, kein freier Parameter -- und genau diese Grenze unterscheidet einen
Kanal für echte Kopplung von einer nachträglichen Zahlenanpassung.

Dieser Kanal ist zugleich der Grund, warum der Fundamentalsektor nicht auf ein starres Skelett
zusammenfällt. Ein vollständig geschlossenes Modell -- jeder Wert ein erzwungener Fixpunkt --
wäre überdeterminiert: keine kontingente Auswahl innerhalb eines $Z_3$-Orbits, keine
Flavour-Mischung, keine physikalische (nicht wegdrehbare) Phase. Gerade die Mischung ist Phase,
nicht Betrag (Dok.~289: geladene Leptonen betrags-dominiert mit kleiner Mischung, Neutrinos
phasen-dominiert mit großer Mischung, der reine Allpass-Sektor). Am schärfsten zeigt sich das an
der Struktur der $CP$-Verletzung: sie verlangt eine irreduzible komplexe Phase -- eine rein reelle,
selbst-adjungierte Struktur ist $CP$-erhaltend --, also genau den Typ des schief-adjungierten
Kanals $\chi=\pi/2$. Der offene Phasenkanal ist damit strukturell von der Art, die $CP$-Verletzung
tragen kann; eine vollständig selbst-adjungierte FFGFT wäre $CP$-gerade und nicht-mischend. Der
Scope bleibt eng: der Kanal trägt Mischung, $CP$ und Kontingenz im \emph{Fundamentalsektor} -- nicht
die makroskopische Vielfalt (Chemie, Strukturbildung), die eigene, davon unabhängige Quellen hat.

\section{Bilanz: die dynamische offene Frage}
Statisch ist die Lage eindeutig: $2/9$ ist unsichtbar im symmetrischen $Q$, blind gegen Betrag
und ganzzahlige Windung, und ein inneres Extremum, das ein freier Parameter herbeizieht
($d\theta^*/d\delta\neq0$). Es ist \emph{nicht} erzwungen wie $c$. Beide Größen teilen dieselbe
flache $T^4$-Arena -- $c$ als Signatur-Rand, $2/9$ als innerer Phasen-Freiheitsgrad --, aber sie
sind verschiedenen Typs.

Damit $2/9$ einen $c$-artigen, erzwungenen Status erhielte, bräuchte es einen \emph{dynamischen}
Fixpunkt, so wie die Eichsymmetrie $m=0$ (und damit $v=c$) erzwingt. Die vorangehenden
Abschnitte lokalisieren diesen Fixpunkt: die \emph{Phase} $\chi=\pi/2$ ist der schief-adjungierte
$i$-Wert, den die Eliminationskette fordert -- der Selektor ist typ-fixiert; offen bleibt allein
der \emph{Betrag} $\delta$ (die golden gedämpfte Oberton-Amplitude). Dessen Feinwert ist dabei
nicht \enquote{noch nicht} hergeleitet, sondern aus der Rekursion allein \emph{beweisbar nicht}
herleitbar: die Nicht-Schließung ist Theorem (einziger Fixpunkt $0$), und die Resonanzdichte der
Durchgangspunkte schließt auch jede rekursionsinterne Auszeichnung eines Wertes aus (Dok.~275,
epistemische Grenze). Die Rekursionsseite ist damit erledigt, keine offene Frage; offen ist allein
die Existenz des äußeren Ankers -- ob eine schief-adjungierte Dynamik gerade $\chi=\pi/2$ trägt
und darüber die Tiefe $\delta$ aus einem inneren Grund \emph{dieser Dynamik} (nicht der Rekursion)
festlegt; gelänge das,
erhielte $2/9$ ein dynamisches, geschütztes Privileg, andernfalls bleibt es kontingent. Die
Eliminationskette (Dok.~282) hat alle statischen Kandidaten -- symmetrisch, zählend-topologisch,
radial, statisch-geometrisch -- ausgeschieden; der einzige Überlebende ist die dynamische
schief-adjungierte Phase, deren Herkunft außerhalb der statischen Geometrie liegt. Genau dort,
nicht im Betrag und nicht in einer flachen Windungszahl, wäre der Wert $2/9$ zu suchen.

\section{Einordnung: Nutzen und Anwendungsnähe}
Die Befunde dieses Dokuments sind Struktur- und Ausschluss-Resultate, keine Technologie. Ihr
unmittelbarer Nutzen liegt bei drei Gruppen. Erstens bei Arbeiten an der Koide-Relation und der
Flavour-Struktur: das No-Go ($2/(9\pi)$ irrational; Betrag und ganzzahlige Windung
counter-blind) schließt topologische und abzählende Herleitungswege für $\theta=2/9$ ab und
erspart deren weitere Verfolgung; die Eliminationskette lokalisiert die verbleibende Suche
präzise im dynamischen schief-adjungierten Phasenkanal. Zweitens methodisch: die
Betrag/Phase-Zerlegung taugt als gemeinsame Audit-Sprache für Framework-Vergleiche
(skript-gedeckte Behauptungen, vorregistrierte Kriterien, matched nulls) und ist vom
FFGFT-Inhalt unabhängig übertragbar. Drittens didaktisch: die Übersetzung zwischen
Teilchenphysik und Nachrichtentechnik (Zirkulant als DFT, Allpass, Injektions-Locking,
Arnold-Zungen) macht den Phasenkanal mit Standard-Ingenieurwerkzeugen rechenbar.

An Messung rührt allein der Phasenkanal selbst: die Strukturaussage, dass Mischung und
$CP$-Verletzung im schief-adjungierten Kanal leben (Dok.~289), hat in den Oszillationsphasen
und $\delta_{CP}$ eine falsifizierbare Form, die laufende und kommende Experimente (DUNE,
Hyper-Kamiokande) vermessen. Eine Anwendung im engeren Sinn -- Gerät, Verfahren, kurzfristiger
Ingenieursnutzen -- folgt daraus nicht: die Werkzeuge laufen von der Nachrichtentechnik in die
Physik, nicht zurück, und solange $\delta$ nicht hergeleitet ist, ist der Ertrag der eines
geschärften offenen Problems, nicht eines Resultats.
